% Shared preamble. Compile main.tex with LuaLaTeX.
\usepackage[letterpaper,margin=1in,headheight=15pt]{geometry}
\usepackage{amsmath,amssymb,amsthm}
\usepackage{fontspec}
\usepackage{unicode-math}
\setmainfont{TeX Gyre Pagella}
\setsansfont{TeX Gyre Heros}[Scale=MatchLowercase]
\setmonofont{TeX Gyre Cursor}[Scale=MatchLowercase]
\setmathfont{TeX Gyre Pagella Math}
\usepackage{microtype}
\usepackage{xcolor}
\definecolor{ReportBlue}{HTML}{17384B}
\definecolor{LinkBlue}{HTML}{176078}
\usepackage{booktabs,longtable,tabularx,array}
\usepackage{enumitem,fancyvrb,fvextra,upquote,caption}
\usepackage{titlesec,fancyhdr,needspace}
\usepackage{xurl}
\usepackage[colorlinks=true,linkcolor=LinkBlue,urlcolor=LinkBlue,citecolor=LinkBlue,
  bookmarksnumbered=true,pdfencoding=auto]{hyperref}
\usepackage{bookmark}
\setlength{\parindent}{0pt}
\setlength{\parskip}{5pt plus 1pt}
\setlength{\emergencystretch}{2em}
\setlength{\tabcolsep}{5pt}
\setlength{\LTpre}{8pt}
\setlength{\LTpost}{8pt}
\renewcommand{\arraystretch}{1.15}
\setlist{topsep=4pt,itemsep=3pt,parsep=1pt,leftmargin=*}
\setcounter{tocdepth}{2}
\setcounter{secnumdepth}{3}
\clubpenalty=10000
\widowpenalty=10000
\displaywidowpenalty=10000
\allowdisplaybreaks[2]
\titleformat{\section}{\Large\bfseries\color{ReportBlue}}{\thesection}{0.7em}{}
\titleformat{\subsection}{\large\bfseries\color{ReportBlue}}{\thesubsection}{0.6em}{}
\titleformat{\subsubsection}{\normalsize\bfseries}{\thesubsubsection}{0.5em}{}
\titlespacing*{\section}{0pt}{18pt}{8pt}
\titlespacing*{\subsection}{0pt}{13pt}{5pt}
\titlespacing*{\subsubsection}{0pt}{10pt}{4pt}
\captionsetup{font=small,labelfont=bf}
\pagestyle{fancy}
\fancyhf{}
\fancyhead[L]{\small\sffamily Distributional AA}
\fancyhead[R]{\small\sffamily November 2026}
\fancyfoot[C]{\small\thepage}
\renewcommand{\headrulewidth}{0.3pt}
\urlstyle{same}
\providecommand{\tightlist}{\setlength{\itemsep}{2pt}\setlength{\parskip}{0pt}}
\providecommand{\pandocbounded}[1]{#1}
\newtheorem{theorem}{Theorem}[section]
\newtheorem{lemma}[theorem]{Lemma}
\newtheorem{proposition}[theorem]{Proposition}
\newtheorem{corollary}[theorem]{Corollary}
\theoremstyle{definition}
\newtheorem{definition}[theorem]{Definition}
\theoremstyle{remark}
\newtheorem{remark}[theorem]{Remark}
\providecommand{\E}{\mathbb E}
\providecommand{\R}{\mathbb R}
\providecommand{\N}{\mathbb N}
\providecommand{\AA}{\mathrm{AA}}
\providecommand{\cAA}{c_{\mathrm{AA}}}
\DeclareMathOperator{\Var}{Var}
\DeclareMathOperator{\Cov}{Cov}
\DeclareMathOperator{\sech}{sech}
\DeclareMathOperator{\tr}{tr}
\DeclareMathOperator{\rank}{rank}
\DeclareMathOperator{\supp}{supp}
\DeclareMathOperator{\diag}{diag}
\pdfstringdefDisableCommands{%
  \def\AA{AA}\def\E{E}\def\Var{Var}\def\Cov{Cov}%
  \def\beta{beta}\def\gamma{gamma}\def\lambda{lambda}%
  \def\mu{mu}\def\rho{rho}\def\le{<=}\def\ge{>=}%
}
